% Compile: pdfLaTeX (or LuaLaTeX). Overleaf-friendly.
% Suspicious spots in hipVS IVF-PQ compute_similarity_kernel for gfx1100 tuning.
\documentclass[aspectratio=169,10pt]{beamer}

\usetheme{Madrid}
\usecolortheme{whale}
\setbeamertemplate{navigation symbols}{}
\setbeamertemplate{footline}[frame number]
\setbeamerfont{title}{size=\Large\bfseries}
\setbeamerfont{frametitle}{size=\large\bfseries}

\usepackage{booktabs}
\usepackage{listings}
\usepackage{xcolor}
\usepackage{hyperref}

\lstset{
  basicstyle=\ttfamily\scriptsize,
  keywordstyle=\color{blue!70!black}\bfseries,
  commentstyle=\color{gray},
  breaklines=true,
  columns=fullflexible,
  frame=single,
  framerule=0.4pt,
  xleftmargin=0.2em,
  xrightmargin=0.2em,
  aboveskip=0.3em,
  belowskip=0.2em,
  showstringspaces=false,
  language=C++
}

\title{hipVS IVF-PQ --- \texttt{compute\_similarity\_kernel}}
\subtitle{Optimization hunt on RDNA3 (\texttt{gfx1100}) \\
\small Source: \texttt{ROCm-DS/hipVS} \texttt{release/rocmds-25.10} \\
\texttt{cpp/src/neighbors/ivf\_pq/ivf\_pq\_compute\_similarity\_impl.cuh}}
\author{Konstantin Kolchin}
\date{13 August 2026}

\begin{document}

%----- title -----
\begin{frame}
  \titlepage
  \vspace{0.2em}
  \begin{center}
    \footnotesize
    Lab: search-only rocprof $\approx$ \textbf{91\%} time in
    \texttt{ivf\_pq::detail::compute\_similarity\_kernel}
    (LUT build + PQ scan). Hints are hypotheses --- measure after each change.
  \end{center}
\end{frame}

%----- map -----
\begin{frame}{Kernel map (what burns time)}
  \footnotesize
  One block $=$ one (\texttt{query}, \texttt{probe}) pair.
  \begin{enumerate}
    \item Optional \textbf{PrecompBaseDiff} into shared memory
    \item Build \textbf{LUT} $[\texttt{pq\_dim}\times 2^{\texttt{pq\_bits}}]$
          (often in LDS / shared)
    \item For each list vector: unpack PQ codes $\rightarrow$ gather LUT
          $\rightarrow$ accumulate score \quad\textbf{(dominant)}
    \item Optional fused local top-$k$ + \texttt{atomicMin} on
          \texttt{query\_kths}
  \end{enumerate}
  \vspace{0.4em}
  \begin{block}{This deck}
    Quote the suspicious spots; one-line hope per spot.
    Not a complete rewrite plan.
  \end{block}
\end{frame}

%----- spot 1: hot score loop -----
\begin{frame}[fragile]{Spot 1 --- per-vector score (hot loop)}
\begin{lstlisting}
score = ivfpq_compute_score<OutT, LutT, vec_t, PqBits>(
    pq_dim,
    reinterpret_cast<...>(pq_thread_data),
    lut_scores,
    early_stop_limit);
\end{lstlisting}
  \vspace{0.3em}
  \begin{alertblock}{Hope (1 line)}
    This call is the PQ scan; any fewer loads / better LUT locality /
    wider vectorization here moves the 91\% needle.
  \end{alertblock}
\end{frame}

%----- spot 2: compute_score body -----
\begin{frame}[fragile]{Spot 2 --- \texttt{ivfpq\_compute\_score} chunk walk}
\begin{lstlisting}
for (; pq_dim >= kChunkSize; pq_dim -= kChunkSize) {
  *pq_codes.vectorized_data() = *pq_head;
  pq_head += kIndexGroupSize;          // interleaved stride
  ivfpq_compute_chunk<...>(score, pq_code, pq_codes,
                           lut_head, lut_end);
  if (score >= early_stop_limit) { return score; }
}
\end{lstlisting}
  \vspace{0.3em}
  \begin{alertblock}{Hope (1 line)}
    Stride \texttt{kIndexGroupSize} and early-exit branches are RDNA3-sensitive;
    retune group size / reduce divergent early-stop.
  \end{alertblock}
\end{frame}

%----- spot 3: bit unpack -----
\begin{frame}[fragile]{Spot 3 --- PQ bit unpack + LUT gather}
\begin{lstlisting}
// inside ivfpq_compute_chunk (recursive template)
uint8_t code = pq_code & kPqMask;
pq_code >>= PqBits;
score += OutT(lut_head[code]);   // gather from LUT
lut_head += kPqShift;            // next subspace book
\end{lstlisting}
  \vspace{0.3em}
  \begin{alertblock}{Hope (1 line)}
    Random LUT gathers + scalar unpack dominate; try denser unrolls,
    half/uint8 LUT, or subspace tiling that hits LDS banks better on RDNA3.
  \end{alertblock}
\end{frame}

%----- spot 4: kIndexGroupSize -----
\begin{frame}[fragile]{Spot 4 --- interleave group $=$ warp assumption}
\begin{lstlisting}
// ivf_pq.hpp (layout contract used by the kernel)
#if defined(CUCO_USE_WARPSIZE_32) && CUCO_USE_WARPSIZE_32
constexpr static uint32_t kIndexGroupSize = 32;
#else
constexpr static uint32_t kIndexGroupSize = raft::warp_size(); // 64 on AMD
#endif
\end{lstlisting}
  \vspace{0.3em}
  \begin{alertblock}{Hope (1 line)}
    Layout/coalescing assumes NVIDIA-style groups; gfx1100 wave32 vs 64
    mismatch wastes bandwidth --- keep packer and scan on the same group size.
  \end{alertblock}
\end{frame}

%----- spot 5: LUT build -----
\begin{frame}[fragile]{Spot 5 --- LUT construction (before the scan)}
\begin{lstlisting}
for (uint32_t i = threadIdx.x; i < lut_size; i += blockDim.x) {
  ...
  do {
    float pq_c = *cur_pq_center;
    cur_pq_center += PqShift;   // strided codebook walk
    switch (metric) {           // L2 / IP / Cosine
      case L2Expanded: { ... score += diff * diff; } break;
      case InnerProduct: { ... score -= q * pq_c; } break;
    }
  } while (++j < j_end);
  lut_scores[i] = LutT(score);
}
\end{lstlisting}
  \vspace{0.3em}
  \begin{alertblock}{Hope (1 line)}
    Metric \texttt{switch} + strided codebook reads inflate LUT build;
    specialize L2-only kernel and vectorize \texttt{pq\_center} loads.
  \end{alertblock}
\end{frame}

%----- spot 6: smem lut -----
\begin{frame}[fragile]{Spot 6 --- LUT in shared vs global}
\begin{lstlisting}
if constexpr (EnableSMemLut) {
  lut_scores = reinterpret_cast<LutT*>(smem_buf);
} else {
  lut_scores += lut_size * blockIdx.x;  // device LUT
}
// later: compute_similarity_select tries
//   (PrecompBaseDiff, EnableSMemLut) =
//   (true,true) -> (false,true) -> (true,false)
\end{lstlisting}
  \vspace{0.3em}
  \begin{alertblock}{Hope (1 line)}
    Heuristics pick shmem/L1 carveout for NVIDIA; re-tune
    \texttt{preferred\_shmem\_carveout} / force variants on gfx1100.
  \end{alertblock}
\end{frame}

%----- spot 7: select occupancy -----
\begin{frame}[fragile]{Spot 7 --- launch heuristics (host)}
\begin{lstlisting}
constexpr double kTargetOccupancy = 0.75;
// grow n_threads_min from warpSize ...
// cudaFuncSetAttribute(...PreferredSharedMemoryCarveout...)
// pick candidate with best occupancy / locality_hint
\end{lstlisting}
  \vspace{0.3em}
  \begin{alertblock}{Hope (1 line)}
    Block size and carveout were not fitted on RDNA3; a gfx1100 grid of
    $(n\_threads, carveout, EnableSMemLut)$ may beat the stock picker.
  \end{alertblock}
\end{frame}

%----- spot 8: precomp -----
\begin{frame}[fragile]{Spot 8 --- \texttt{PrecompBaseDiff} into LDS}
\begin{lstlisting}
if constexpr (PrecompBaseDiff) {
  // L2: lut_end[i] = query[i] - cluster_center[i];
  // IP:  store float2(query, query*center)
  __syncthreads();
}
\end{lstlisting}
  \vspace{0.3em}
  \begin{alertblock}{Hope (1 line)}
    Extra LDS pressure can steal L1 on RDNA3; for some
    $(\texttt{pq\_dim},\texttt{dim})$ the \texttt{false} variant may win.
  \end{alertblock}
\end{frame}

%----- spot 9: fences -----
\begin{frame}[fragile]{Spot 9 --- barriers before the scan}
\begin{lstlisting}
__threadfence_block();
__syncthreads();
local_topk_t block_topk(topk, lut_end, query_kth);

for (uint32_t i = threadIdx.x; i < n_samples_aligned;
     i += blockDim.x, pq_thread_data += pq_line_width) {
  ...
}
\end{lstlisting}
  \vspace{0.3em}
  \begin{alertblock}{Hope (1 line)}
    Double sync after LUT may be heavier than needed on HIP; test dropping
    \texttt{\_\_threadfence\_block} if LDS visibility is already covered.
  \end{alertblock}
\end{frame}

%----- spot 10: filter in hot path -----
\begin{frame}[fragile]{Spot 10 --- filter check inside sample loop}
\begin{lstlisting}
bool valid = i < n_samples;
if (valid && sample_filter(queries_offset + query_ix, label, i)) {
  score = ivfpq_compute_score<...>(...);
  if (metric == CosineExpanded) { score = OutT(1) + score; }
}
\end{lstlisting}
  \vspace{0.3em}
  \begin{alertblock}{Hope (1 line)}
    Even \texttt{none\_sample\_filter} still sits in the hot path; compile a
    filter-free specialization for the unfiltered Milvus case.
  \end{alertblock}
\end{frame}

%----- spot 11: fused topk -----
\begin{frame}[fragile]{Spot 11 --- fused top-$k$ + atomic}
\begin{lstlisting}
if constexpr (kManageLocalTopK) {
  block_topk.add(score, sample_offset + i);
  ...
  block_topk.done(smem_buf);
  block_topk.store(out_scores, out_indices);
  if (threadIdx.x == 0) {
    atomicMin(query_kths + query_ix, float(out_scores[topk - 1]));
  }
}
\end{lstlisting}
  \vspace{0.3em}
  \begin{alertblock}{Hope (1 line)}
    Warp-sort + cross-probe \texttt{atomicMin} contend on AMD; try
    Capacity$=0$ (score then select) for large batches on gfx1100.
  \end{alertblock}
\end{frame}

%----- spot 12: early stop -----
\begin{frame}[fragile]{Spot 12 --- early stop vs full scan}
\begin{lstlisting}
OutT early_stop_limit = kDummy;
switch (metric) {
  case L2Expanded: case L2SqrtExpanded:
    early_stop_limit = query_kth; break;
  default: break;
}
// inside score loop:
if (score >= early_stop_limit) { return score; }
\end{lstlisting}
  \vspace{0.3em}
  \begin{alertblock}{Hope (1 line)}
    Early-stop saves FLOPs but adds divergence; A/B with
    \texttt{kDummy} always (full scan) on wave32 RDNA3.
  \end{alertblock}
\end{frame}

%----- priority -----
\begin{frame}{Suggested attack order}
  \footnotesize
  \begin{enumerate}
    \item \textbf{Spots 1--4} --- PQ scan + interleave/\texttt{kIndexGroupSize}
          (where 91\% lives)\\
          {\scriptsize First patch in harness:
          \texttt{patches/hipvs/0004-\ldots compute-score-pipeline-pq8\ldots}
          --- software-pipeline loads + \texttt{pq\_bits=8} byte LUT path.
          Apply: \texttt{scripts/apply\_hipvs\_ivf\_pq\_compute\_score\_opt.sh}}
    \item \textbf{Spot 6--7} --- force launch variants / carveout grid on gfx1100
    \item \textbf{Spot 5} --- L2-specialized LUT build
    \item \textbf{Spots 8--12} --- cheaper wins (sync, filter-free, top-$k$ mode,
          early-stop A/B)
  \end{enumerate}
  \vspace{0.5em}
  \begin{block}{Measure}
    Keep search-only rocprof (time-filter after build). Report QPS +
    \% time in \texttt{compute\_similarity\_kernel} per change.
  \end{block}
\end{frame}

\end{document}
